% --- LaTeX Homework Template - S. Venkatraman ---

% --- Set document class and font size ---

\documentclass[letterpaper, 11pt]{article}

% --- Package imports ---

\usepackage{
  amsmath, amsthm, amssymb, mathtools, dsfont,	  % Math typesetting
  graphicx, wrapfig, subfig, float,                  % Figures and graphics formatting
  listings, color, inconsolata, pythonhighlight,     % Code formatting
  fancyhdr, sectsty, hyperref, enumerate, enumitem } % Headers/footers, section fonts, links, lists

% --- Page layout settings ---

% Set page margins
\usepackage[left=0.7in, right=0.7in, bottom=1in, top=1.1in, headsep=0.2in]{geometry}

% Anchor footnotes to the bottom of the page
\usepackage[bottom]{footmisc}



% Set line spacing
\renewcommand{\baselinestretch}{1.2}

% Set spacing between paragraphs
\setlength{\parskip}{1.5mm}

% Allow multi-line equations to break onto the next page
\allowdisplaybreaks

% Enumerated lists: make numbers flush left, with parentheses around them
\setlist[enumerate]{wide=0pt, leftmargin=21pt, labelwidth=0pt, align=left}
\setenumerate[1]{label={(\arabic*)}}

% --- Page formatting settings ---

% Set link colors for labeled items (blue) and citations (red)
\hypersetup{colorlinks=true, linkcolor=blue, citecolor=red}

% Make reference section title font smaller
\renewcommand{\refname}{\large\bf{References}}

% --- Settings for printing computer code ---

% Define colors for green text (comments), grey text (line numbers),
% and green frame around code
\definecolor{greenText}{rgb}{0.5, 0.7, 0.5}
\definecolor{greyText}{rgb}{0.5, 0.5, 0.5}
\definecolor{codeFrame}{rgb}{0.5, 0.7, 0.5}

% Define code settings
\lstdefinestyle{code} {
  frame=single, rulecolor=\color{codeFrame},            % Include a green frame around the code
  numbers=left,                                         % Include line numbers
  numbersep=8pt,                                        % Add space between line numbers and frame
  numberstyle=\tiny\color{greyText},                    % Line number font size (tiny) and color (grey)
  commentstyle=\color{greenText},                       % Put comments in green text
  basicstyle=\linespread{1.1}\ttfamily\footnotesize,    % Set code line spacing
  keywordstyle=\ttfamily\footnotesize,                  % No special formatting for keywords
  showstringspaces=false,                               % No marks for spaces
  xleftmargin=1.95em,                                   % Align code frame with main text
  framexleftmargin=1.6em,                               % Extend frame left margin to include line numbers
  breaklines=true,                                      % Wrap long lines of code
  postbreak=\mbox{\textcolor{greenText}{$\hookrightarrow$}\space} % Mark wrapped lines with an arrow
}

% Set all code listings to be styled with the above settings
\lstset{style=code}

% --- Math/Statistics commands ---

% Add a reference number to a single line of a multi-line equation
% Usage: "\numberthis\label{labelNameHere}" in an align or gather environment
\newcommand\numberthis{\addtocounter{equation}{1}\tag{\theequation}}

% Shortcut for bold text in math mode, e.g. $\b{X}$
\let\b\mathbf

% Shortcut for bold Greek letters, e.g. $\bg{\beta}$
\let\bg\boldsymbol

% Shortcut for calligraphic script, e.g. %\mc{M}$
\let\mc\mathcal

% \mathscr{(letter here)} is sometimes used to denote vector spaces
\usepackage[mathscr]{euscript}

% Convergence: right arrow with optional text on top
% E.g. $\converge[w]$ for weak convergence
\newcommand{\converge}[1][]{\xrightarrow{#1}}

% Normal distribution: arguments are the mean and variance
% E.g. $\normal{\mu}{\sigma}$
\newcommand{\normal}[2]{\mathcal{N}\left(#1,#2\right)}

% Uniform distribution: arguments are the left and right endpoints
% E.g. $\unif{0}{1}$
\newcommand{\unif}[2]{\text{Uniform}(#1,#2)}

% Independent and identically distributed random variables
% E.g. $ X_1,...,X_n \iid \normal{0}{1}$
\newcommand{\iid}{\stackrel{\smash{\text{iid}}}{\sim}}

% Equality: equals sign with optional text on top
% E.g. $X \equals[d] Y$ for equality in distribution
\newcommand{\equals}[1][]{\stackrel{\smash{#1}}{=}}

% Math mode symbols for common sets and spaces. Example usage: $\R$
\newcommand{\R}{\mathbb{R}}   % Real numbers
\newcommand{\C}{\mathbb{C}}   % Complex numbers
\newcommand{\Q}{\mathbb{Q}}   % Rational numbers
\newcommand{\Z}{\mathbb{Z}}   % Integers
\newcommand{\N}{\mathbb{N}}   % Natural numbers
\newcommand{\F}{\mathcal{F}}  % Calligraphic F for a sigma algebra
\newcommand{\El}{\mathcal{L}} % Calligraphic L, e.g. for L^p spaces

% Math mode symbols for probability
\newcommand{\pr}{\mathbb{P}}    % Probability measure
\newcommand{\E}{\mathbb{E}}     % Expectation, e.g. $\E(X)$
\newcommand{\var}{\text{Var}}   % Variance, e.g. $\var(X)$
\newcommand{\cov}{\text{Cov}}   % Covariance, e.g. $\cov(X,Y)$
\newcommand{\corr}{\text{Corr}} % Correlation, e.g. $\corr(X,Y)$
\newcommand{\B}{\mathcal{B}}    % Borel sigma-algebra

% Other miscellaneous symbols
\newcommand{\tth}{\text{th}}	% Non-italicized 'th', e.g. $n^\tth$
\newcommand{\Oh}{\mathcal{O}}	% Big-O notation, e.g. $\O(n)$
\newcommand{\1}{\mathds{1}}	% Indicator function, e.g. $\1_A$

% Additional commands for math mode
\DeclareMathOperator*{\argmax}{argmax}    % Argmax, e.g. $\argmax_{x\in[0,1]} f(x)$
\DeclareMathOperator*{\argmin}{argmin}    % Argmin, e.g. $\argmin_{x\in[0,1]} f(x)$
\DeclareMathOperator*{\spann}{Span}       % Span, e.g. $\spann\{X_1,...,X_n\}$
\DeclareMathOperator*{\bias}{Bias}        % Bias, e.g. $\bias(\hat\theta)$
\DeclareMathOperator*{\ran}{ran}          % Range of an operator, e.g. $\ran(T) 
\DeclareMathOperator*{\dv}{d\!}           % Non-italicized 'with respect to', e.g. $\int f(x) \dv x$
\DeclareMathOperator*{\diag}{diag}        % Diagonal of a matrix, e.g. $\diag(M)$
\DeclareMathOperator*{\trace}{trace}      % Trace of a matrix, e.g. $\trace(M)$

% Numbered theorem, lemma, etc. settings - e.g., a definition, lemma, and theorem appearing in that 
% order in Section 2 will be numbered Definition 2.1, Lemma 2.2, Theorem 2.3. 
% Example usage: \begin{theorem}[Name of theorem] Theorem statement \end{theorem}
\theoremstyle{definition}
\newtheorem{theorem}{Theorem}[section]
\newtheorem{proposition}[theorem]{Proposition}
\newtheorem{lemma}[theorem]{Lemma}
\newtheorem{corollary}[theorem]{Corollary}
\newtheorem{definition}[theorem]{Definition}
\newtheorem{example}[theorem]{Example}
\newtheorem{remark}[theorem]{Remark}

% Un-numbered theorem, lemma, etc. settings
% Example usage: \begin{lemma*}[Name of lemma] Lemma statement \end{lemma*}
\newtheorem*{theorem*}{Theorem}
\newtheorem*{proposition*}{Proposition}
\newtheorem*{lemma*}{Lemma}
\newtheorem*{corollary*}{Corollary}
\newtheorem*{definition*}{Definition}
\newtheorem*{example*}{Example}
\newtheorem*{remark*}{Remark}
\newtheorem*{claim}{Claim}
\newcommand{\Tau}{\mathcal{T}}
\newcommand{\PSet}{\mathcal{P}}


% --- Left/right header text (to appear on every page) ---

% Include a line underneath the header, no footer line
\pagestyle{fancy}
\renewcommand{\footrulewidth}{0pt}
\renewcommand{\headrulewidth}{0.4pt}

% Left header text: course name/assignment number
\lhead{MAM4000W, Analysis II -- Assignment 1}

% Right header text: your name
\rhead{Alex Cristaudo, CRSALE010}

% --- Document starts here ---

\begin{document}
\subsection*{Question 1}
\begin{itemize}
    \item \textbf{Positivity:} We can always cover $E$ with itself, and choose some $i \in I$ and let $ J = \{i\}$ and $E_i = E$, then $E = \bigcup \limits _{j \in J} E_j$, and thus $\mu (E) \ge \mu _i (E) \ge 0$ since $\mu _i$ is a positive measure.     \item \textbf{Empty Set:} Clearly, as in Positivity, this supremum set is not empty so suppose $(E_i)_{i \in J} \subset \mathcal{E}$ is pairwise disjoint and $\emptyset = \bigcup \limits _{i \in J} E_i$ and $J \subset I$ is at most countable. Since $\emptyset = \bigcup \limits _{i \in J} E_i$, we must have $E_i = \emptyset$ for all $i \in J$. So \[
            \sum \limits _{i \in J} \mu _i (E_i) = \sum \limits _{i\in J} \underbrace{\mu _i (\emptyset)}_{=0} = 0
        \] since each $\mu _i$ is a positive measure. Thus $\mu (\emptyset) = 0$
    \item $\mathbf{\sigma}$\textbf{-additivity:} Suppose that $(E_n)_{n\in \mathbb{N}} \subset \mathcal{E}$ is pairwise disjoint. Let $\varepsilon > 0$. Then, for each $n\in \mathbb{N}$, we have a family $(E_{nm})_{m \in J_n} \subset \mathcal{E}$ that is pairwise disjoint and $E_n = \bigcup \limits _{m \in J_n} E_{nm}$ and $J_m \subset I$ is at most countable, and \[\mu (E_n) - \varepsilon 2^{-n} < \sum \limits_{m \in J_n} (\mu _m(E_{nm}))  \] so \[
    	\sum \limits _{n \in \mathbb{N}} (\mu (E_n)) - \varepsilon < \sum \limits _{n\in \mathbb{N}}  \sum \limits _{m \in J_n} (\mu _m (E_{nm})) 
    \]
    
    
    Since the summands are nonnegative, we can use Tonelli's Theorem for sums and swap indices. To do so, let $J = \bigcup \limits _{n\in \mathbb{N}} J_n$,  $N_j = \{n \in \mathbb{N} : j \in J_n\}$ for $j \in J$ and $A_j := \bigcup \limits _{n \in N_j} E_{nj}$. $J$ is at most countable, as a countable union of at most countable sets. Since $(E_{nm})_{n \in \mathbb{N}, m \in J_n}$ is pairwise disjoint, so is $(A_j)_{j \in J}$. Moreover, \[\bigcup \limits _{n\in \mathbb{N}}E_n = \bigcup \limits _{n\in \mathbb{N}, j \in J_n} E_{nm} = \bigcup \limits _{j \in J} A_j\] 
 \[ \text{So }
    \mu \left ( \bigcup \limits _{n \in \mathbb{N}} E_n \right ) \ge \sum \limits _{j \in J}\mu _j (A_j) = \sum \limits _{j \in J} \mu _j \left ( \bigcup \limits _{n \in N_j} E_{nj} \right)= \sum \limits _{j \in J} \sum \limits _{n \in N_j} \mu _j(E_{nj)} = \sum \limits _{n \in \mathbb{N}} \sum \limits _{j \in J_n} \mu _j(E_{nj)}> \sum \limits _{n \in \mathbb{N}} (\mu (E_n)) - \varepsilon 
    \]
    
    This is true for all $\varepsilon > 0$ and thus\[
     \mu \left ( \bigcup \limits _{n \in \mathbb{N}} E_n \right ) \ge \sum \limits _{n \in \mathbb{N}} (\mu (E_n)) 
    \]
To show the other inequality, let $E = \bigcup \limits _{n \in \mathbb{N}} E_n = \bigcup \limits _{j \in J} B_j$ for $(B_j)_{j \in J} \subset \mathcal{E}$ pairwise disjoint and $J \subset I$ is at most countable. Now, noticing that $(B_j \cap E_n)_{j \in J}$ is p.d. and  \[
\sum \limits _{ j\in J} \mu _j(B_j) = \sum \limits _{ j\in J} \mu _j \left (\bigcup \limits _{n\in \mathbb{N}}B_j \cap E_n \right) = \sum \limits _{ j\in J} \sum  _{n \in \mathbb{N}} \mu _j(B_j \cap E_n) = \sum \limits _{n \in \mathbb{N}} \sum \limits _{ j \in J} \mu _j(B_j \cap E_n) 
\]Then, since $E_n = \bigcup \limits _{j \in J} (B_j \cap E_n)$, with $(B_j \cap E_n)_{j \in J}$ being pairwise disjoint, we get \[\sum \limits _{j \in J} \mu _j (B_j \cap E_n) \le \mu (E_n)\]
So finally,
\[
\sum \limits _{j \in J} \mu _j (B_j) \le  \sum \limits _{n \in \mathbb{N}} \mu (E_n)
\]
meaning
 \[
\mu \left ( \bigcup \limits _{n \in \mathbb{N}} E_n \right ) \le \sum \limits _{n \in \mathbb{N}} (\mu (E_n)) 
\]
making $\mu$ $\sigma$-additive and thus a positive measure
    
    \end{itemize}
\subsection*{Question 2}

\begin{enumerate}[label=(\alph*)]
\item {\begin{itemize}
\item \textbf{Empty set:} For any set $A \in \mathcal{B}(X)$, clearly $\mu (A) \ge 0$ so $\Phi (E) \ge 0$ for all $E \subset X$ (note that $X$ itself is open, so $\{\mu (A): A \text{ is open}, A \supset E\}$ is nonempty). On the other hand, $\emptyset$ is open and $\mu (\emptyset) = 0$, so $\Phi (\emptyset) = 0$

\item $\mathbf{\sigma}$\textbf{-subadditivity:} Let $A  \subset X$, $(A_n)_{n \in \mathbb{N}} \subset \mathcal{P}(X)$ such that $A \subset \bigcup _{n \in \mathbb{N}}A_n$. Let $\varepsilon > 0$. For each $n \in \mathbb{N}$, there exists an open set $O_n \supset A_n$ such that, since $\Phi$ is an infimum, and that $\mu$ is finite, \[\Phi (A_n) + \varepsilon 2^{-n} > \mu (O_n)\] Then $\bigcup \limits _{n \in \mathbb{N}} O_n \supset \bigcup \limits _{n\in \mathbb{N}} A_n \supset A$ and $\bigcup \limits _{n \in \mathbb{N}} O_n$ is open. So, since $\mu$ is $\sigma$-subadditive, \[
\Phi  ( A ) \le \mu \left ( \bigcup \limits _{n\in \mathbb{N}} O_n\right ) \le \sum \limits _{n \in \mathbb{N}} \mu (O_n) \le \sum \limits _{n \in \mathbb{N}} (\Phi (A_n)) + \varepsilon
\]  
This is true for all $\varepsilon > 0$ and so
\[
\Phi (A) \le  \sum \limits _{n \in \mathbb{N}} \Phi (A_n)
\]


\end{itemize}
}


\item {Let $E_1, E_2 \subset X$ with $d(E_1, E_2) = r > 0$. From $\sigma$-subadditivity, $\Phi(E_1 \cup E_2) \le \Phi (E_1) + \Phi (E_2)$. Now $d(E_1, E_2) =  \inf \limits _{(e_1, e_2) \in E_1 \times E_2} d(e_1, e_2)$. For each $e_1 \in E_1$, clearly $B(e_1; \frac{1}{2}r) \cap E_2 = \emptyset$. Consider the open sets $U_1 = \bigcup \limits _{e_1 \in E_1} B(e_1; \frac{1}{2}r)$ and $U_2 = \bigcup \limits _{e_2 \in E_2} B(e_2; \frac{1}{2}r)$. Then $E_1 \cup E_2 \subset U_1 \cup U_2$ and $U_1 \cap U_2 = \emptyset$ (as if $x \in U_1 \cap U_2$, then there exist some $e_1 \in E_1, e_2 \in E_2$ such that $d(e_1,x) < \frac{1}{2}r$ and $d(e_2, x) < \frac{1}{2}r$, then $d(e_1, e_2) < r$, a contradiction).  Let $O \supset E_1 \cup E_2$ be open. Then $O \cap U_1 \supset E_1$ and $O \cap U_2 \supset E_2$ and both are open, so, using the additivity of the measure $\mu$,
\[
\mu (O) \ge \mu (O \cap (U_1 \cup U_2)) = \mu ((O \cap U_1) \cup (O \cap U_2)) = \mu (O \cap U_1) + \mu (O \cap U_2) \ge \Phi (E_1) + \Phi (E_2)
\]
So $\Phi (E_1 \cup E_2 ) \ge \Phi (E_1) + \Phi (E_2)$ and
\[
\Phi(E_1) + \Phi (E_2) = \Phi(E_1 \cup E_2)
\]
}
\end{enumerate}


\subsection*{Question 3}
$(\Rightarrow )$ Suppose $f$ is measurable. $f^2$ is measurable since $f^2 = f\cdot f$, as the product of measurable functions, and $\{x \in X: f(x) > 0\} = f^{-1}((0, \infty )) \in \mathcal{E}$ since $(0, \infty)$ is a Borel set. \\
$(\Leftarrow)$ Let $t \in \mathbb{R}$. We need to show that $\{x \in X: f(x) > t\} \in \mathcal{E}$. The case $t=0$ is given.\\ If $t  > 0$, then we have that $f(x)>t \iff f(x) > 0 \text{ and } |f(x)| > |t| \iff f(x) > 0 \text{ and } f(x)^2 > t^2 $, giving \\ $\{x \in X: f(x) > t\} = \{x \in X: f(x) > 0\} \cap \{x \in X: f(x)^2 > t^2\} ) \in \mathcal{E}$. \\ If $ t < 0$, then we have that $f(x) > t \iff f(x) > 0 \text{ or } (f(x) \le 0 \text{ and } |f(x)| < |t|) \iff f(x) > 0 \text{ or } (f(x) \le 0 \text{ and } f(x)^2 < t^2)  $, giving \\$\{x \in X:f(x) > t\} = \{x \in X: f(x) > 0\} \cup (\{x \in X: f(x) \le 0\} \cap \{x \in X: f(x)^2  < t^2\}) \in \mathcal{E}$ \\ So for all $t\in \mathbb{R}$, $f^{-1}((t, \infty)) \in \mathcal{E}$ making $f$ is measurable


\end{document}
